\section{Evaluation}
\label{sec:eval}

A demonstration paper must show that the system works, not re-establish the
methodology (done in \citealp{dynamicmcpbench2026}). We therefore evaluate
three properties of the \emph{studio}: that its verdicts match the batch
pipeline, that it is fast enough to drive live, and that visitors can reach
the verdict flip unaided.

\subsection{E1: Studio-vs-batch verdict agreement}

The studio is only credible if it reproduces the pipeline's verdicts. We
score the same set of \emph{(task, candidate)} pairs two ways---through the
studio back-end and through the original batch pipeline---and measure
agreement on the Tier-1 effect verdict, which is deterministic and so should
agree exactly. Table~\ref{tab:agreement} reports the result.

% E1 -- studio-vs-batch verdict agreement. Numbers from
% dmcp-studio/experiments/e1_agreement.py (report:
% docs/experiments/studio-e1-agreement.md). Tier-1 is deterministic; the
% pre-registered bar is exact 100% agreement.
\begin{table}[t]
  \centering \small
  \begin{tabular}{lcc}
    \toprule
    & \textbf{comparisons} & \textbf{agreement} \\
    \midrule
    Overall verdict        & 118 & \textbf{100.0\%} \\
    Per-checkpoint (Tier-1) & 708 & \textbf{100.0\%} \\
    \bottomrule
  \end{tabular}
  \caption{\textbf{E1: studio-vs-batch agreement.} Each of 118 deterministic
    \emph{(spec, candidate)} pairs (22 pass / 96 fail) is scored through the
    studio backend and through the original batch pipeline (the \texttt{dmcp
    eval} CLI, a separate process). The studio's JSON round-trip and result
    shaping reproduce the pipeline's deterministic Tier-1 verdict exactly
    --- overall and for every checkpoint --- meeting the pre-registered 100\%
    bar. Tier-2 (the non-deterministic LLM equivalence judge) is off the demo's
    critical path.}
  \label{tab:agreement}
\end{table}


\noindent The pre-registered target is 100\% Tier-1 agreement; any
discrepancy indicates the adapter diverged from the pipeline and is a bug to
fix, not a tolerance to accept. Across 118 deterministic pairs (22 pass, 96
fail) the studio reproduces the batch verdict exactly---both overall and for
all 708 per-checkpoint comparisons---so the demo shows the pipeline's own
scorer, not a re-implementation. Tier-2 (the LLM effect-equivalence judge) is
reported separately, since it is not deterministic and is off the demo's
critical path.

\subsection{E2: Per-stage latency budget}

Whether a stage runs LIVE or from a REPLAY fixture in the booth is decided
by its wall-clock cost. We measure the per-stage \emph{compute} of the REPLAY
booth path through the real routes (inter-event pacing disabled, so we time
the system rather than the cosmetic animation delay).
Table~\ref{tab:latency} reports the budget.

% E2 -- per-stage latency budget.
% TODO(E2): fill from experiments/e2_latency.py. Do NOT fabricate; leave TBD
% until the run is committed. Columns: median wall-clock, #tool calls, mode.
\begin{table}[t]
  \centering \small
  \begin{tabular}{lccc}
    \toprule
    \textbf{stage} & \textbf{median} & \textbf{calls} & \textbf{booth mode} \\
    \midrule
    goal generation & \textit{TBD} & --- & \textit{TBD} \\
    explore (live)   & \textit{TBD} & TBD & \textit{TBD} \\
    distill          & \textit{TBD} & --- & \textit{TBD} \\
    score            & \textit{TBD} & TBD & \textit{TBD} \\
    \bottomrule
  \end{tabular}
  \caption{\textbf{E2: per-stage latency budget.} Median wall-clock and
    tool-call counts per stage in LIVE mode, with the resulting booth
    decision (run LIVE vs.\ serve a REPLAY fixture). \emph{Numbers pending
    the E2 run.}}
  \label{tab:latency}
\end{table}


\noindent Every stage is sub-millisecond and cold-start to first verdict is
${\approx}1.3$\,ms---roughly four orders of magnitude under the 30\,s target
(\S\ref{sec:limitations})---so the only delay a visitor perceives is the
deliberate, configurable inter-event pacing ($\sim$0.45\,s/call). The decision
rule follows: the booth runs REPLAY for every stage. LIVE exploration, which
would be dominated by LLM and network seconds, is offered as an explicit proof
run rather than the graded path---consistent with scoring under deterministic
replay throughout.

\subsection{E3: Showcase fixtures}

We curate and freeze the three verdict-flip fixtures of \S\ref{sec:system},
plus the worked-example default (the AAPL/MSFT/GOOGL financial-comparison task):
each is a candidate trajectory scored by the \emph{real} Tier-1 evaluator (E1),
so the demo deterministically reproduces each disagreement with the pipeline's
own verdicts rather than a re-implementation. Each fixture carries a one-line
rationale and is served by the REPLAY back-end; LIVE re-derives the same TaskSpec
against the read-only servers.

\subsection{E5: Informal usability}

We ran an informal pass with a handful of users unfamiliar with the system,
measuring time-to-first-verdict-flip and noting points of confusion; no
formal user study is claimed. The observations fed directly back into the
interface (labelling of the toggle, the equivalence-set hint, and the
LIVE/REPLAY affordance). \emph{[Fill from the E5 session: median
time-to-first-flip and the one or two fixes folded back.]}
